\chapter{Rotator Cuff Tendinopathy / Impingement Syndrome}
\index{Rotator Cuff Tendinopathy / Impingement Syndrome}
\label{sec:Rotator Cuff Tendinopathy / Impingement Syndrome}

\section{Overview}
Rotator cuff tendinopathy encompasses a spectrum of disorders affecting the tendons of the rotator cuff muscles (supraspinatus, infraspinatus, teres minor, subscapularis), ranging from tendinosis to partial- and full-thickness tears. Impingement syndrome refers to mechanical compression of the rotator cuff tendons and subacromial bursa between the humeral head and the acromion. Both conditions commonly cause shoulder pain and functional limitation.

\section{Classification}
\begin{description}[font=\normalfont\bfseries,leftmargin=3em,labelwidth=2.5em]
  \item[Cause] Overuse, tendon degeneration, or mechanical compression under the acromion
  \item[Organ System] Bones/Musculoskeletal
  \item[Pathophysiology] Repetitive overhead activity leads to collagen fiber microtrauma, failed healing response with mucoid degeneration, neovascularization, and altered tenocyte function. Impingement involves narrowing of the subacromial space, with bursal inflammation and mechanical abrasion of the tendon against the coracoacromial arch.
  \item[Duration] Chronic
  \item[Onset] Gradual (overuse) or acute (after injury)
  \item[Mode of Transmission] Not applicable
  \item[Clinical Course] Insidious onset of shoulder pain; may progress to tendon tear if unmanaged.
  \item[Severity] Mild to moderate; can become severe with tendon tearing
  \item[Extent of Disease] Localized to the shoulder
  \item[Age Group] More common in adults aged 40--60 years; younger athletes also affected
  \item[Epidemiology] Affects 15--30\% of the general population; one of the most common causes of shoulder pain. High prevalence in overhead athletes and occupations involving repetitive arm elevation.
  \item[ICD-11 Code] FA50.0 (rotator cuff tendinopathy), FA50.1 (impingement syndrome)
\end{description}

\section{Causes}
Rotator cuff tendinopathy results from repetitive microtrauma due to overhead activities, tendon degeneration with aging, or acute injury. Impingement syndrome is caused by anatomical narrowing of the subacromial space, often due to acromial morphology (hook-shaped acromion), subacromial spurs, or thickening of the coracoacromial ligament.

\section{Risk Factors}
\begin{itemize}[noitemsep]
  \item Age over 40 years
  \item Repetitive overhead activities (throwing, swimming, painting, construction)
  \item Poor shoulder biomechanics or scapular dyskinesis
  \item Acromial morphology (type III hooked acromion)
  \item Rotator cuff muscle weakness or imbalance
  \item Prior shoulder injury or instability
\end{itemize}

\section{Signs and Symptoms}
\begin{itemize}[noitemsep]
  \item Deep, aching shoulder pain, especially with overhead activity
  \item Night pain that disrupts sleep
  \item Pain with reaching behind the back
  \item Weakness or difficulty lifting the arm
  \item Catching or grinding sensation with shoulder movement
  \item Difficulty performing daily activities
\end{itemize}

\section{Progression and Stages}
Tendinopathy progresses through reactive tendinopathy (reversible), tendinosis disrepair (partially reversible), and degenerative tendinopathy (irreversible structural changes). Impingement follows Neer's staging: stage I (edema, reversible), stage II (fibrosis, tendinitis), and stage III (tendon tear). Gradual worsening of pain and weakness may culminate in a full-thickness rotator cuff tear.

\section{Complications}
\begin{itemize}[noitemsep]
  \item Progression to partial- or full-thickness rotator cuff tear
  \item Adhesive capsulitis (frozen shoulder) secondary to immobility
  \item Chronic shoulder pain and disability
  \item Persistent loss of shoulder range of motion
  \item Reduced quality of life
  \item Emotional and psychological effects
\end{itemize}

\section{Diagnosis}
\begin{itemize}[noitemsep]
  \item Physical examination (Hawkins--Kennedy test, Neer impingement sign, empty can test)
  \item Assessment of shoulder range of motion and strength
  \item Lab tests (not routinely needed)
  \item Imaging tests (ultrasound or MRI for tendon assessment; plain radiographs for acromial morphology)
\end{itemize}

\section{Prevention}
\begin{itemize}[noitemsep]
  \item Strengthen rotator cuff and scapular stabilizer muscles
  \item Use proper technique for overhead activities and sports
  \item Avoid sudden increases in training load or intensity
  \item Regular check-ups
\end{itemize}

\section{Treatment Options}
\begin{itemize}[noitemsep]
  \item Physical therapy focused on rotator cuff strengthening and scapular stabilization
  \item Nonsteroidal anti-inflammatory drugs (NSAIDs) for pain and inflammation
  \item Subacromial corticosteroid injection for short-term relief
  \item Activity modification and ergonomic adjustments
  \item Surgical options: arthroscopic subacromial decompression, rotator cuff repair
  \item Lifestyle changes
\end{itemize}

\section{Living with It}
\begin{itemize}[noitemsep]
  \item Follow treatment plan
  \item Regular appointments
  \item Adjust daily habits
  \item Support groups
  \item Keep learning
\end{itemize}

\section{Outlook}
Most patients with rotator cuff tendinopathy and impingement improve with conservative management within 6--12 weeks. Approximately 70--80\% respond to non-operative treatment. For those requiring surgery, outcomes are generally favorable with appropriate rehabilitation. Chronic cases or large tendon tears have a less favorable prognosis, with persistent weakness or pain in a subset of patients.
